% Transcription — fonds Alexandre Grothendieck, shelfmark 118, batch 2 (pages 21-40).
% Established from the facsimile archives/batches/118/batch-02.pdf (a mirror of
% https://grothendieck.umontpellier.fr/118.pdf; no cover sheet, PDF page k =
% folder page 20+k), per the skill transcribe-grothendieck, with the 700 dpi
% tiles of archives/tiles/118/ and 400-600 dpi re-crops; diagonal margins
% re-rotated with magick (35-75 degrees).
%
% Pass: Opus 5.5 (claude-opus-5-5), 2026-09-24 — first pass, unchecked against the pages by a human.
%
% Nineteen of the twenty pages are transcribed. Skipped: page 37, a blank
% sheet (only the archivists' pencil number). There is no typescript and no
% text in another hand in the batch; nothing is summarised. His own
% pagination « 3 » ... « 7 » stands at the top of pages 21-25 (recorded in a
% \note on each), continuing batch 1's run; from page 26 on no number of his
% is visible. Page 39 is a landscape double sheet written in two columns.
%
% What the batch is about. It continues batch 1 on the localisateur
% fondamental W of (Cat) and the W_A it induces on a presheaf category A^.
% Pages 21-22: a cocartesian (and cartesian) square of monomorphisms in A^
% gives, via A_{/X}, a square in Cat of open immersions to which the
% earlier statement applies; the Corollary (his margin: « Thm ») that
% pushouts of such squares preserve W; then « Mais attention »: in Cat a
% cocartesian square along a closed immersion need not transport W, with a
% counter-example built from the cone, suspension and Θ of a groupoid Z with
% π_1 ≠ 1. Pages 23-25: Hot_A as a category of fractions of the homotopy
% category of A^ by the homotopy classes of W_A, and the claim that these
% admit a calculus of left fractions; the axioms (i)-(iv) on a class 𝔐 of
% arrows (factorisation through a homotopism, cocartesian stability,
% separating intervals, finite sums), with his margins that A must be a
% local test category and that (iii) is useless for A^; the use of Loc. 3
% for finite sums. Pages 26-28: Corollary 1 — Hot(W) has sums and (Cat) →
% Hot(W), A^ → Hot_A commute with them; a Lemma for a homotopy interval
% structure (finite sums, f ~ g iff f_j ~ g_j); the NB that (Cat) → \bar Cat
% does not commute with infinite sums (the infinite zigzag I_∞), and his
% guess that Hot(W) nevertheless does. Pages 29-30: the bijectivity of
% Hom_Hot(∐X_j, Y) → ∏Hom_Hot(X_j, Y) — surjectivity via a common
% amalgamated Y', a Lemma on a family of weak equivalences Y → Y'_j, and the
% reduction of injectivity to a Kan-type condition. Pages 31-36: Corollary
% 2 — finite products in Hot(W), preserved by (Cat) → Hot(W) and A^ →
% Hot(W) for A a strict test category (his margin: totally W-aspherical
% suffices); surjectivity and injectivity; a Lemma listing equivalent
% conditions (i) A totally W-aspherical, (ii)/(ii bis) W_A stable under
% finite products / under - × id_Z, (iii), (iv), (iv bis) the functors
% commute with finite products, proved on pages 33-34 and summed up in a
% scheme of implications; page 35's NB that (iv) is stable under sums while
% (i) is not, the groupoid counter-example, and the question of infinite
% products of weak equivalences, which page 36 sets aside as less important
% than infinite sums. Pages 38-39: a Theorem — a map of categories filtered
% by closed (resp. open) subcategories indexed by an ordered set I, which is
% a weak equivalence on each stage, is a weak equivalence — via the
% functor φ: X → I, φ(x) = Inf{i | x ∈ X_i}, and L 3; a Corollary in A^,
% and « Cor ? » on unions of W-aspherical X_i over a connected I. Pages
% 39-40: Thm « L'ens. W est saturé » and Cor « W_A est saturé » — a first
% proof, cancelled by two diagonal strokes, with f, g, f' and α, β ∈ W; then
% on page 40 the working proof: f with γ(f) invertible, the calculus of
% fractions gives g_1 with g_1 f_0 ∈ W, a sequence of open immersions
% X_0 → X_1 → X_2 → … whose consecutive composites lie in W, X_∞ its
% colimit, and i_0 as colimit of the X_0 → X_{2n} — « donc on gagne !! ».
%
% Notation fixed for the batch. \underline{W} for his underlined W (the
% localisateur of Cat) and \underline{W}_A; plain W where he does not
% underline; \overline{A^{\wedge}}, \overline{\underline{W}_A},
% \overline{\mathrm{Cat}} for his barred objects (homotopy categories);
% A^{\wedge} for his « A^ », \widehat{A} where he draws a hat (pages 34, 35);
% \mathfrak{M} for the round script M of page 24; \mathrm{Hot}_A,
% \mathrm{Hot}(\underline{W}); \gamma for the localisation functor;
% (\mathrm{Cat}) with or without parentheses as written. Underlined words
% are \emph. His « OPS » (on peut supposer), « ssi », « tt », « p.ex. »,
% « équiv. », « imm. », « asph. » are kept; his « = » for « même(s) » is
% noted where it occurs (pages 22, 29, 30). The abbreviation read « tot. »
% (totalement W-asphérique) on pages 31-35 looks like « LA »/« tA » and is
% marked \uncertain throughout.
%
% Notable uncertainties: the second half of page 21 (the sign read ∫ before
% A_{/X'}); page 22's diagonal margin (« 1-équiv. … groupoïdes
% fondamentaux ») and the indices u'_0, w_0 of the example; page 24's two
% diagonal margins; page 25's left margin, nearly vertical and mostly \ill;
% page 26's margin on « finies inutile »; page 29's line a), crowded and
% partly \ill; page 38's header line, not read, and its oblique margin on
% Inf; page 39's cancelled first proof and the small sketch at its foot;
% page 40's bracketed « [Il suffit] bien sûr … ». Cross-references: « Loc. 3 »
% (page 25) and « L 3 » (page 39) refer to conditions/lemmas numbered in
% batch 1 or elsewhere; transcribed as written.
\documentclass[11pt,a4paper]{article}
% Le préambule commun à toutes les transcriptions du fonds Grothendieck.
%
% Il tient en une idée : **l'incertitude se compose, elle ne se résout pas.**
% Une transcription qui lisse un mot douteux a détruit la seule chose pour
% laquelle on la faisait — un lecteur doit pouvoir savoir, à chaque endroit,
% ce qui était sur le feuillet et ce qui a été ajouté. Les six macros qui
% suivent sont donc l'appareil critique, et non de la décoration.
%
% Les mêmes commandes sont reconnues par `scripts/render.mjs`, qui produit la
% vue de lecture du site. Une macro ajoutée ici doit l'être aussi là-bas,
% faute de quoi le rendu échouera bruyamment — ce qui est voulu : un rendu
% silencieusement faux est pire qu'un rendu absent.

\ProvidesPackage{grothendieck}[2026/08/08 Transcriptions du fonds Grothendieck]

\usepackage{amsmath,amssymb,amsthm}
\usepackage{mathrsfs}
\usepackage{stmaryrd}
% Les diagrammes commutatifs sont partout dans ce fonds : c'est la langue
% même de la géométrie algébrique de Grothendieck, et une transcription qui
% les rendrait en prose ne transcrirait rien.
\usepackage{tikz-cd}
% Français par défaut — c'est la langue des feuillets — mais l'anglais est
% chargé aussi : la traduction et le résumé sont des documents anglais, et la
% typographie française (espaces avant les deux-points) y serait une faute.
% Ces documents basculent par \selectlanguage{english} en tête de corps.
\usepackage[english,french]{babel}
\usepackage{fontspec}
\usepackage{geometry}
\geometry{a4paper,margin=25mm,marginparwidth=28mm}
\usepackage{marginnote}
\usepackage{xcolor}
% ulem, not soul. `soul`'s \st and \ul are built on a character-by-character
% scan that XeLaTeX's font machinery breaks: the document fails with an
% undefined control sequence rather than degrading. ulem does the same two jobs
% and is Unicode-engine safe. `normalem` stops it hijacking \emph, which the
% transcriptions use for Grothendieck's own emphasis.
\usepackage[normalem]{ulem}
\usepackage{hyperref}
\hypersetup{colorlinks=true,linkcolor=black,urlcolor=black,pdfborder={0 0 0}}

% --- Métadonnées du lot -----------------------------------------------------
% Reprises telles quelles par le script de rendu, qui les affiche en tête de
% la vue de lecture. Elles fixent ce que le fichier prétend couvrir : sans
% elles, une transcription flotte sans point d'attache dans un fonds de seize
% mille feuillets.

\newcommand{\folder}[1]{\gdef\grothFolder{#1}}
\newcommand{\batch}[1]{\gdef\grothBatch{#1}}
\newcommand{\leaves}[2]{\gdef\grothFirst{#1}\gdef\grothLast{#2}}
\newcommand{\foldertitle}[1]{\gdef\grothFoldertitle{#1}}
\gdef\grothFolder{?}\gdef\grothBatch{?}\gdef\grothFirst{?}\gdef\grothLast{?}\gdef\grothFoldertitle{}

% --- Le résumé --------------------------------------------------------------
%
% Il ouvre la lecture modernisée, sous le titre « Résumé », et il est composé
% en petit. Ce n'est pas un ornement : c'est le seul passage du document qui
% ne soit pas des mathématiques, et un lecteur doit pouvoir le reconnaître
% avant de l'avoir lu — puis le sauter s'il connaît déjà le sujet, ou s'y
% arrêter s'il ne le connaît pas. Le filet qui le suit marque la reprise du
% corps.

\newenvironment{resume}
  {\small\setlength{\parskip}{0.45em}}
  {\par\medskip\hrule\medskip}

% --- Mots-clés ---------------------------------------------------------------
%
% En anglais, en clôture du résumé de la lecture modernisée : le vocabulaire
% moderne sous lequel chercher ce que les feuillets construisent. Le manifeste
% du site (`npm run manifest`) les extrait pour étiqueter la cote entière —
% cette ligne est la source unique des tags, il n'y a pas de fichier à côté.
\newcommand{\keywords}[1]{\par\smallskip\noindent
  {\footnotesize\textsc{Keywords} --- \textit{#1}}\par}

% --- L'appareil critique ----------------------------------------------------

% Le numéro de feuillet, en marge. C'est l'ancre qui permet de régler un
% désaccord sur une formule en regardant *un* feuillet plutôt qu'une chemise
% de six cents. Le site s'en sert aussi pour tourner les pages du fac-similé
% quand on fait défiler la transcription.
\newcommand{\leaf}[1]{%
  \marginnote{\footnotesize\color{gray!70}\textsf{#1}}%
  \label{leaf:#1}\ignorespaces}

% Illisible : on ne devine pas. Un mot inventé qui se lit comme les autres est
% le pire résultat possible.
\newcommand{\ill}{\textcolor{red!60!black}{[\,\ldots\,]}}

% Lecture douteuse : le mot est donné, et signalé comme incertain.
\newcommand{\uncertain}[1]{\uline{#1}}

% Ajout de l'éditeur — une lettre manquante, un mot sous-entendu.
\newcommand{\add}[1]{\textcolor{blue!50!black}{[#1]}}

% Biffé par Grothendieck lui-même. Ce qu'il a rayé fait partie du document :
% une rature dit souvent où la pensée a changé de direction.
\newcommand{\struck}[1]{\sout{#1}}

% Note du transcripteur, sur l'état matériel du feuillet.
\newcommand{\note}[1]{\par\smallskip{\small\color{gray!60!black}\textit{[#1]}}\par\smallskip}

% Note marginale de Grothendieck, distincte du corps du texte.
\newcommand{\marginal}[1]{\par\smallskip\noindent\hspace{1em}%
  {\small\color{gray!60!black}#1}\par\smallskip}

% --- Titre ------------------------------------------------------------------

\AtBeginDocument{%
  \begin{center}
    {\large\bfseries Fonds Alexandre Grothendieck --- cote n\textsuperscript{o}~\grothFolder}\\[2pt]
    {\itshape\grothFoldertitle}\\[4pt]
    {\small feuillets \grothFirst--\grothLast\ (lot \grothBatch)}\\[2pt]
    {\footnotesize\color{gray!60!black}Transcription établie d'après le fac-similé numérisé
     par l'Université de Montpellier.\\
     Les lectures incertaines sont soulignées, les illisibles notées [\,\ldots\,],
     les ajouts entre crochets.}
  \end{center}
  \vspace{1em}}

\endinput


\folder{118}
\batch{2}
\pages{21}{40}
\dating{[à partir de 1983]}
\watermark{Édition de démonstration}
\foldertitle{Cohomology properties of maps in (Cat) : notes manuscrites (s.d.).}

\begin{document}

\page{21}
\note{sa p. 3}

Soit maintenant

\begin{tikzcd}
Y \arrow[r, hook] \arrow[d] & X \arrow[d] \\
Y' \arrow[r, hook] & X'
\end{tikzcd}

diagramme cocartésien dans un $\widehat{A}$, avec $Y \hookrightarrow X$
\emph{\uncertain{et} aussi cartésien}\note{ajouté dans l'interligne,
souligné} (donc $Y' \hookrightarrow X'$ mono.). Considérons le diagramme
\struck{cartésien} correspondant dans Cat

\begin{tikzcd}
A_{/Y} \arrow[r, hook] \arrow[d] & A_{/X} \arrow[d] \\
A_{/Y'} \arrow[r, hook] & A_{/X'}
\end{tikzcd}

il est cocartésien (car $i_A : A^{\wedge} \to (\mathrm{Cat})$ commute aux
$\varinjlim$) et cartésien (car $i_A$ commute aux produits fibrés) et de
plus les flèches \ill{} sont fibrantes, de plus $A_{/Y} \to A_{/X}$ et
$A_{/Y'} \to A_{/X'}$ \struck{sont} \ill{} des mono.\ fibrantes donc des
immersions ouvertes. On est donc dans la situation déjà énoncée, donc
$\int A_{/X'}$\note{le signe qui précède $A_{/X'}$ a la forme d'un
$\int$ ; lecture incertaine} \uncertain{réalise}
$\int \begin{pmatrix} A_{/Y} \to A_{/X} \\ \downarrow \\ A_{/Y'} \end{pmatrix}$,

\marginal{\textbf{Thm}}\note{doublement souligné, en marge gauche, en face
du corollaire}
\emph{Corollaire} Si $(Y \to X) \in \underline{W}$ (resp.\ $Y \to Y'$
\struck{\ill{}} $\in W$) alors $(Y' \to X') \in \underline{W}$\note{$Y'
\to X'$ est récrit sur une première écriture} (resp.\ $X \to X'
\in \underline{W}$).

\page{22}
\note{sa p. 4}

\emph{Mais attention} Il n'est pas vrai \emph{dans Cat} que si on a un
diagramme cocartésien

\begin{tikzcd}
Y \arrow[r, hook, "i"] \arrow[d, "g"'] & X \arrow[d, "f"] \\
Y' \arrow[r, hook, "i'"'] & X'
\end{tikzcd}

si $i \in \underline{W}$, a-t-on $i' \in \underline{W}$ ? \ill{} $Y, X$
asphériques (\emph{Ex} $Y' = \Delta_0$) \textbf{Non}\note{ces trois lignes
sont écrites à droite du carré ; « Non » est doublement souligné}

$i$ immersion fermée (ou ouverte), alors $[\, i \in \underline{W}
\Rightarrow i' \in \underline{W}$, ni que$\,]$ $g \in \underline{W}
\Rightarrow f \in \underline{W}$, et même\note{il écrit « = » pour
« même », ici et deux mots plus loin} pour l'équiv.\ faible habituelle, et
même si $g$ est \struck{\ill{}} cofibrante \struck{\ill{}}, et $i$
\ill{} immersion fermée, $Y' \simeq \Delta_0$\note{la fin de la phrase,
depuis « si $g$ », est écrite dans l'interligne}.

\marginal{Il semble \uncertain{vrai} que si $g$ \struck{\ill{}}
\uncertain{1-équiv.}, i.e.\ induit une \struck{\ill{}} équiv.\ sur les
\uncertain{groupoïdes fondamentaux}, \ill{} \uncertain{$f$ aussi}.}\note{note
écrite en diagonale dans la marge gauche, à hauteur du carré ; « cocart. »
est ajouté au-dessus de la première ligne}
$[\;]$ ?

\emph{Ex} Considérons une $Z$ dans (Cat), et

\begin{tikzcd}
Z \arrow[r, hook] \arrow[d, hook, "\alpha"'] & \check{C}(Z)\ [\supset e_{u'_0}] \arrow[d] \\
Y = C(Z) \arrow[r, hook, "\text{imm. f. (surj.)}"'] \arrow[d, "g \in W"'] & \Sigma(Z) = X\ [\supset e_{w_0}] \arrow[d, "f"] \\
Y' = e \arrow[r, hook, "\text{imm. f.}"'] & \Theta(Y) = X'\ [\supset e_{u'_0}]
\end{tikzcd}

\note{les deux carrés sont marqués « cocart » ; la flèche $\alpha$ porte
« immersion fermée » ; la flèche horizontale du haut porte un petit signe
illisible ; les indices $u'_0$, $w_0$ des crochets sont lus sans
certitude, et dans le premier crochet un signe est biffé devant
$e_{u'_0}$}

Ici $g \in W$, et $X' \simeq \Delta^1$ ssi \struck{\ill{}} $Z$
\struck{\ill{}} est 0-connexe. (En général, $\mathrm{Hom}(e_{u'_0},
e_{y'_0}) \simeq \pi_0(Z)$\note{indices incertains}). Mais quand il en est
ainsi, $f \in \underline{W}$ ssi $\Sigma(Z)$ est $\underline{W}$-asph.,
ce qui n'est pas toujours le cas (prendre p.ex.\ $Z$ un groupoïde réduit à
un objet, avec $\pi_1 \neq 1$).

\page{23}
\note{sa p. 5, chiffre récrit}

\marginal{Il faut supposer que $A$ soit catégorie test locale.}\note{en
diagonale dans la marge gauche, séparé de la suivante par un trait}
\marginal{itou dans \struck{\ill{}} Cat}

\note{devant « Corollaire », un « 4 » et un numéro biffé}
\emph{Corollaire} Considérons $\overline{A^{\wedge}}$ la catégorie
localisée de $A^{\wedge}$ par les homotopismes (relatifs :
$h_{\underline{W}_A}$), de sorte que $\mathrm{Hot}_A$ apparaît comme une
catégorie de fractions
\[
  \mathrm{Hot}_A \simeq \overline{\underline{W}_A}^{\,-1}\, \overline{A^{\wedge}}
\]
Je dis que $\overline{\underline{W}_A}$ permet un calcul des fractions à
gauche. \struck{\ill{}} Soit en effet $Y \xrightarrow{i} X$,
$g \in \overline{\underline{W}_A}$ : $Y \to Y'$\note{les
$Y$, $Y'$ de cette ligne sont surchargés} dans $\overline{A^{\wedge}}$,
\ill{} compléter en

\begin{tikzcd}
Y \arrow[r, "i"] \arrow[d, "g"'] & X \arrow[d, "f"] \\
Y' \arrow[r, "i'"'] & X'
\end{tikzcd}

\ill{} $f \in \overline{\underline{W}_A}$.

On va factoriser \struck{\ill{}} $i$ en \struck{mono} $Y \xrightarrow{i_1}
X_1 \xrightarrow{h} X$, \struck{$Y \to (X \times I)$} (cf.\ plus bas)
\ill{} $h$ homotopisme, on est ramené au cas où $i$ \struck{\ill{}} mono,
et on prend $X'$ somme amalgamée. D'autre part, on doit montrer que

\page{24}
\note{sa p. 6}

\[
  X' \xrightarrow{\;s\;} X \overset{f}{\underset{g}{\rightrightarrows}} Y
  \quad \text{de } \overline{A^{\wedge}}
\]
avec $s \in \overline{\underline{W}_A}$, \ill{} \ill{} dans
\[
  X \overset{f}{\underset{g}{\rightrightarrows}} Y \xrightarrow{\;t\;} Y'
\]
avec $t \in \overline{\underline{W}_A}$\note{la lettre, ici et dans
l'étiquette de $Y \to Y'$, est lue $t$ d'après le diagramme qui suit ; elle
pourrait être un $r$}.
\textbf{OPS} $s$ mono.

\begin{tikzcd}[column sep=small]
 & I \times X \arrow[rr, dashed] & & Y' \\
I \times X' \arrow[r, hook, "{[\in \underline{W}_A]}"'] & M \arrow[u, "{[\in \underline{W}_A]}"] \arrow[rr] & & Y \arrow[u, dashed, "t"'] \\
X' \sqcup X' \arrow[u] \arrow[r, hook, "\in \underline{W}_A"'] & X \sqcup X \arrow[u] \arrow[urr] & &
\end{tikzcd}

\note{au-dessus de $I \times X'$, une source biffée ($X'$ surchargé) envoie
une flèche marquée $\in \underline{W}_A$ vers $I \times X$ ; la flèche
$X' \sqcup X' \to I \times X'$ est tracée en gras}

\marginal{L'\uncertain{énoncé} des conditions (i) à (iv) \uncertain{inspiré}
par le cas de (Cat) et non celui de $A^{\wedge}$, où (iii) est inutile,
pourvu que dans (ii) il suffise \ill{} \ill{} $\in \underline{W}$.}\note{en
diagonale dans la marge droite, une flèche renvoyant à « les choses
suivantes »}
On va utiliser les choses suivantes \struck{de $(I_\alpha)$, la famille
d'int.} :
$\exists$ une classe de flèches $\mathfrak{M} \subset
\mathrm{Fl}(\underline{M})$\note{$\mathfrak{M}$ rend sa lettre ronde, écrite
sur une autre biffée} (ici $\underline{M} = A^{\wedge}$) \struck{\ill{}}, avec
\begin{itemize}
  \item[(i)] toute $f : Y \to X$ se factorise en $Y \xrightarrow{i} X_1
  \xrightarrow{h} X$, avec $i \in \mathfrak{M}$, $h$ homotopisme
  \struck{rel.\ $(I_\alpha)$}
  \item[(ii)] Si \struck{\ill{}}

\begin{tikzcd}
Y \arrow[r, "i"] \arrow[d, "f"'] & X \arrow[d, "g"] \\
Y' \arrow[r] & X'
\end{tikzcd}

  cocartésien, avec $i \in \mathfrak{M}$, et $f \in \underline{W}$, alors
  $g \in W$
  \item[(iii)] $\exists$ famille génératrice d'intervalles $I_\alpha$,
  telle que pour tout $X$, $X \sqcup X \to X \times I_\alpha$ soit
  $\in \mathfrak{M}$\marginal{ce sont donc des intervalles séparants}
  \item[(iv)] \emph{Sommes finies} \struck{\ill{}} \uncertain{existent}
  dans $\underline{M}$, et \struck{\ill{} \ill{}} $f, g \in W \Rightarrow
  f \sqcup g \in W$.
\end{itemize}
\marginal{\uncertain{modifier} \ill{}}\note{dans la marge, à hauteur de
(ii)--(iii), deux lignes barrées d'un trait courbe}

\marginal{Dans toute $(M, W)$ \struck{\ill{}} \uncertain{munie d'un}
$\mathfrak{M}$ (compatible) $W$ [\uncertain{var.} les homotopismes sont
dans $W$] \ill{} (i) \ill{}, $W^{-1}M = \overline{W}^{-1}\overline{M}$, et
$\overline{W}$ admet un calcul des fractions à gauche, \uncertain{si}
\ill{} les conditions (i) à (iv).}\note{en diagonale dans la marge
gauche, en face des conditions (i)--(iv)}

\page{25}
\note{sa p. 7 ; en haut à droite, un mot biffé, peut-être
\struck{\uncertain{Fukuda}}}

Dans cet argument, il a fallu \uncertain{moralement}\note{ajouté au-dessus
de la ligne} utiliser un intervalle $I$ \emph{séparant} (cf condition
(iii)). Donc dans le cas $M = A^{\wedge}$, il faut disposer d'un intervalle
$I$ \emph{séparant} dans $A^{\wedge}$ \struck{donc $A$} qui de plus
définisse une \struck{\ill{}} \uncertain{structure}\note{récrit au-dessus du
mot biffé} homotopique sur $A^{\wedge}$ \struck{plus fine \ill{}}
compatible avec $W$ -- et pour ceci, il faut supposer que $I \to e$ est
$\underline{W}_A$-asphérique [ou du moins supposer que $\delta_0$ et
$\delta_1 : e \to I$ \struck{\ill{}} \emph{\ill{}}\note{mot doublement
souligné} \struck{les deux ext.} \struck{\ill{}} même image dans
$\underline{W}$ ....] -- donc, pratiquement, que \struck{\ill{}} $A$ soit
une catégorie-test \emph{locale}.
\marginal{\ill{} inutile, cf.\ \ill{} \ill{} si \ill{} \struck{\ill{}}
catégorie $A$}\note{marge gauche, écrit presque verticalement, en face des
lignes 2 à 8 ; une note plus bas, à hauteur de « catégorie-test locale »,
est biffée et ne se lit pas}
\marginal{\emph{Ah}}

On a utilisé aussi la propriété (pour card $J = 2$) que pour une famille de
flèches $(f_j)_{j \in J}$ dans $A^{\wedge}$, $f_j \in \underline{W}_A$ pour
tout $j \in J$ $\Rightarrow$ $\coprod_J f_j \in \underline{W}_A$. On est
ramené à voir la même dans (Cat), pour $\underline{W}$. Cela résulte alors
de Loc.\ 3 (condition de localisation) appliquée à

\begin{tikzcd}[column sep=small]
X \arrow[rr] \arrow[dr] & & Y \arrow[dl] \\
 & J &
\end{tikzcd}

\struck{Coroll.}

\page{26}

\emph{Corollaire 1} Dans $\mathrm{Hot} = \mathrm{Hot}(\underline{W})$, les
sommes \struck{\ill{}} \uncertain{arbitraires} existent, et les foncteurs
\uncertain{canoniques}
\[
  (\mathrm{Cat}) \to \mathrm{Hot}(\underline{W}), \qquad
  A^{\wedge} \to \mathrm{Hot}_A \to \mathrm{Hot}(\underline{W})
\]
y commutent (on se ramène au cas de (Cat), \uncertain{sachant} que
$i_A : A^{\wedge} \to (\mathrm{Cat})$ commute aux sommes), i.e.\ les
\emph{\ill{}} dans \ill{} \ill{}.
\marginal{\emph{finies} inutile, \ill{} \uncertain{quelconque}, \ill{}
\uncertain{plus} \ill{}}\note{marge gauche, en oblique, à hauteur des trois
premières lignes ; une flèche en part vers $A^{\wedge}$, une autre vers
« sans contradiction » écrit au-dessus de la ligne suivante, lecture
incertaine}

\emph{Démonstration} Soit $(M, W)$ avec $W \subset \mathrm{Fl}(M)$,
$W$ \struck{\ill{}} permettant un calcul des fractions à droite.
\struck{Si dans} Alors \struck{$M$ les} le foncteur canonique $M \to W^{-1}M$
commute aux \emph{lim} finies, \struck{\ill{}} et en particulier aux
sommes finies.\note{les quatre lignes, de « $W$ permettant » à « sommes
finies », sont marquées d'un trait vertical en marge}

On applique ceci à $(\overline{\mathrm{Cat}}, \overline{\underline{W}})$,
il faut voir que les sommes finies y existent (\ill{} la \ill{} \ill{}
\ill{} \ill{} quelque chose !), et que $\mathrm{Cat} \to
\overline{\mathrm{Cat}}$ y commute. \struck{Dans} Mais \ill{} de
\uncertain{façon} générale

\page{27}

\marginal{\emph{NB} $M$ \uncertain{stable} par sommes finies
\uncertain{suffit}}\note{en oblique dans le coin supérieur gauche}
\emph{Lemme} \struck{Soit} Soit $(M, h)$ structure homotopique par
intervalles (homotopy interval structure), avec $M$ stable par sommes
$\{$finies$\}$, et les sommes $\{$finies$\}$ étant universelles. Alors
$M \to W_h^{-1}M$ commute aux sommes finies.

Cela revient à ceci : pour
\[
  \begin{matrix} f = (f_j)_{j \in J} \\ g = (g_j)_{j \in J} \end{matrix}
  \; : \; \underbrace{\textstyle\coprod_J X_j}_{X}
  \overset{f}{\underset{g}{\rightrightarrows}} Y
\]
$f \underset{h}{\sim} g$ ssi $f_j \underset{h}{\sim} g_j$ $\forall j \in J$.
\note{ces lignes sont marquées d'un trait vertical en marge}

\emph{Dém} $\Rightarrow$ trivial, pour prouver $\Leftarrow$ on note qu'il
existe \emph{un \uncertain{même} intervalle} $I$ tel que
$f_j \underset{h_I}{\sim} g_j$ \struck{\ill{}} (car $J$ est fini),
\struck{on prend} comme
\[
  I \times X \simeq \coprod_J I \times X_j
\]
on trouve une homotopie composée \ill{}

\begin{tikzcd}
I_N \times X \arrow[r] & X \\
I_N \times X_j \arrow[u, hook] \arrow[r] & X_j \arrow[u, hook]
\end{tikzcd}

\note{au-dessus de $I_N \times X$, un premier facteur biffé ; les
inclusions sont notées « $\cup$ » sur la page}
(avec $N$ assez grand pour convenir pour tous les $j$ …). \uncertain{Ça
gagne} !

\page{28}

\emph{NB} Il semble \emph{faux} que $(\mathrm{Cat}) \to
\overline{\mathrm{Cat}}$ commute aux sommes infinies -- il suffit de
prendre des familles
\[
  \begin{matrix} f_n \\ g_n \end{matrix} \Big\} \; X_n \rightrightarrows Y
  \qquad (n \in \mathbb{N})
\]
telles que $f_n \simeq g_n$, mais qu'il faille une suite de \uncertain{plus en
plus} \ill{} d'homotopies élém.\ pour les joindre. (prendre p.ex.\
$X_n = \Delta_0$, $f_n$, $g_n$ donnés par des pts de $Y$, prendre
\[
  Y = I_\infty = (\;\underbrace{\to\leftarrow}\,\underbrace{\to\leftarrow}\;\cdots\;)
\]
somme amalgamée d'une infinité de copies de $X_n$, \struck{Mais pour
\ill{}} la somme\note{« la somme » est écrit au-dessus des mots biffés}
$\coprod X_n$ \emph{n'est pas} une somme dans $\overline{\mathrm{Cat}}$,
\uncertain{car}\note{la parenthèse ouverte après « prendre p.ex. » n'est pas
refermée ; le signe $\coprod$ devant $X_n$ paraît barré}

Je suppose pourtant que dans $\mathrm{Hot}(\underline{W})$, les sommes
quelconques existent et que $\mathrm{Cat} \to \mathrm{Hot}(\underline{W})$
y commute ??

\page{29}

\marginal{\struck{Correspondance avec \ill{}}}\note{en haut à droite,
biffé de traits en boucle ; le dernier mot, sans doute un nom, ne se lit
pas}
\struck{Pourtant, soit $X$}

Il reviendrait au même de dire que pour $(X_j)_{j \in J}$ et $Y$ dans Cat,
si $X = \coprod X_j$,
\[
  \mathrm{Hom}_{\mathrm{Hot}}(X, Y) \to \prod_{j \in J} \mathrm{Hom}_{\mathrm{Hot}}(X_j, Y)
\]
\note{devant $X$ et $Y$ dans le premier terme, deux signes noircis}
est bijectif. Mais avec le calcul des fractions, ni l'\struck{bijecti}injectivité
ni la surjectivité n'est claire, elles auraient l'air plutôt fausses, pour
les mêmes\note{il écrit « = » pour « même(s) »} raisons. C'est
\uncertain{gênant} = bien gênant ! \ill{}

a) Il se pourrait que pour \struck{deux} \struck{\ill{}} \ill{} des
\uncertain{sections} \ill{} $(\varphi_j)_{j \in J}$\note{avant la famille,
un signe « = » suivi d'une lettre, peut-être son abréviation de « mêmes »}, il n'existe pas de
$(Y \to Y') \in W$ \emph{commun pour tous les $\varphi_j$}, tels que
$X_j \overset{\varphi_j}{\dashrightarrow} Y$ se représente par
$X_j \xrightarrow{f_j} Y'$. A priori, on trouve $Y \hookrightarrow Y'_j$
(\uncertain{immersion} \ill{} \ill{}), mais \uncertain{peut-être} pas
\uncertain{pour la} \ill{}. \ill{} \uncertain{proj.} : les \ill{},
\ill{} pour partie finie

\page{30}

$J_0$ de $J$ \uncertain{fini}, $Y'_{J_0} = \coprod_{j \in J_0,\, Y} Y'_j$,
et $Y'_J = \varinjlim Y'_{J_0}$, on a encore $Y \to Y'$ $\in
\underline{W}$. Donc il y a quand même\note{« = » pour « même »}
surjectivité ! Mais soient maintenant $f, g : X
\overset{f}{\underset{g}{\rightrightarrows}} Y'$ (sous $Y$, avec
$Y \hookrightarrow Y'$)\note{il écrit $Y$ au-dessus de $Y'$, relié par une
flèche verticale} qui \struck{donnent} sont tels que $\gamma(f_j) =
\gamma(g_j)$ $\forall j \in J$, a-t-on $\gamma(f) = \gamma(g)$ ? Mais oui,
par le même argument -- conduisant au lemme :
\marginal{\ill{} \ill{} !}\note{en oblique dans la marge gauche, souligné de
deux traits}

\emph{Lemme} Soient $Y \xrightarrow{\alpha_j} Y'_j$ \struck{des} une
famille d'équivalences faibles, alors il existe $Y
\overset{\alpha}{\hookrightarrow} Y'$ une équivalence faible \struck{telle
que} (immersion ouverte ou fermée si les $\alpha_j$ le sont)
\struck{telles que} et des $\beta_j : Y'_j \to Y'$ telles que $\beta_j
\alpha_j = \alpha$. (De plus, si les $\alpha_j$ sont des immersions
ouvertes ou imm.\ fermées, on peut prendre pour $Y'$ la somme amalgamée des
$Y'_j$)\note{au-dessus de « Y' » dans la parenthèse du lemme, « $Y
\hookrightarrow Y'$ » est ajouté}

\begin{tikzcd}
Y \arrow[r, hook, "\overline{\alpha}_j"] \arrow[dr, hook, bend right=30, "\alpha"'] & \overline{Y}'_j \arrow[r, "i_j"] \arrow[d, "\text{comm.}"'] & Y'_j \arrow[l, dashed, bend left=40, "s_j"] \arrow[dl, dashed, "\beta_j"] \\
 & Y' &
\end{tikzcd}

\[
  s_j i_j = \mathrm{id}_{\overline{Y}'_j}
\]
\note{schéma en marge gauche ; « comm. » est écrit deux fois, dans les deux
triangles}

On est ramené à prouver ceci : soient $f, g : X = \coprod X_j \to Y$,
tels que $f_j \sim g_j$ $\forall j \in J$, alors $\gamma(f) = \gamma(g)$.
(Mais attention on n'a pas en général $f \sim g$ !) \emph{Il faut une
\uncertain{variation} \ill{} des types « \uncertain{Kan} » sur $Y$}, ou
plutôt sur un $Y'$, où $(Y \to Y') \in \underline{W}$

\page{31}

\emph{Corollaire 2} \emph{Dans} $\mathrm{Hot}(\underline{W})$, les
produits finis existent, et les foncteurs canoniques
\[
  (\mathrm{Cat}) \to \mathrm{Hot}(\underline{W}), \qquad
  A^{\wedge} \to \mathrm{Hot}(\underline{W}) \quad (A \text{ catégorie test
  stricte})
\]
y commutent.
\marginal{il suffit même que $A$ soit \uncertain{tot.}\ $\underline{W}$-asph.}\note{ajout
relié par une flèche à « $A$ catégorie test stricte » ; lecture
incertaine de la fin}

\marginal{\struck{\ill{} \ill{}} Mais c'est clair dans (Cat) si on traite
directement}\note{en oblique dans la marge gauche, à hauteur des quatre
lignes suivantes}
On est ramené à prouver que $A^{\wedge} \to [\mathrm{Hot}(\underline{W})
\simeq]\ \underline{W}_A^{-1} A^{\wedge}$ commute aux produits finis, i.e.\
que
\[
  \mathrm{Hom}_{\underline{W}_A^{-1}A^{\wedge}}\Bigl(X, \prod Y_j\Bigr)
  \to \prod \mathrm{Hom}_{\underline{W}_A^{-1}A^{\wedge}}(X, Y_j)
\]
\note{le premier $\prod Y_j$ est surchargé}
est bijectif, pour toute famille finie $(Y_j)_{j \in J}$ d'objets en
$A^{\wedge}$, et $X$ dans $A^{\wedge}$.

\emph{Surjectivité} \uncertain{Un élément} \ill{} \ill{} \uncertain{est}
provient d'une famille
\[
  f_j : X \to Y'_j \qquad \text{où } Y_j \to Y'_j \text{ équiv.\ faible.}
\]
\struck{Donc} Considérons donc
\[
  f = (f_j) : X \to \prod Y'_j = Y'
\]
il définit $X \to Y$ qui convient, \emph{si on sait que} $\prod Y_j \to
\prod Y'_j$ est $\in \underline{W}$.

\page{32}

\marginal{via localisation forte [si $A$ \struck{\ill{}} \ill{}
$\underline{W}$-asph.]}\note{écrit verticalement le long du bord gauche,
avec deux flèches doubles montantes ; entre la marge et les conditions,
des flèches $\Downarrow$ et $\Updownarrow$ relient (i) à (ii bis) et
(ii bis) à (iii)}
C'est le cas si $A$ catégorie-test stricte, plus généralement si elle est
\uncertain{tot.}\ $\underline{W}$-asphérique.\note{l'abréviation, ici, en
p.~31 et p.~33, ressemble à « LA » ou « tA » ; lue « tot. » (totalement)
d'après la condition $X \times a \to X \in \underline{W}_A$ de p.~33, sans
certitude}

\emph{Lemme} (?) Soit $A$ dans (Cat) \struck{catégorie \ill{}}.
Conditions équiv.\ ($\Rightarrow$)
\begin{itemize}
  \item[(i)] $A$ \struck{est} \uncertain{tot.}\ $\underline{W}$-asphérique,
  \struck{il \ill{}}
  \item[(ii)] \struck{Produits finis de flèches $\in \underline{W}_A$ sont
  $\in \underline{W}_A$}
  \item[(ii bis)] \struck{(iii)} Si $f : X \to Y$ \struck{dans $A$} est
  dans $\underline{W}_A$ et $Z$ dans $A^{\wedge}$, $f \times \mathrm{id}_Z :
  X \times Z \to Y \times Z$ est dans $\underline{W}_A$
  \item[(iii)] Le foncteur $A^{\wedge} \to (\mathrm{Hot}_A)$
  \struck{$\mathrm{Hot}(\underline{W})$} commute aux produits finis
  \item[(iv)] le foncteur $A^{\wedge} \to (\mathrm{Hot}_A(\underline{W}))$
  commute aux produits finis
\end{itemize}
\note{sous « (iii) », un mot souligné et biffé, peut-être
\struck{\uncertain{Dém}}}

On se ramène à ce lemme \struck{tantôt}. Pour le moment, on suppose
vérifié (ii), qui assure surjectivité de $\ast$.

\emph{Injectivité} ?

\begin{tikzcd}
 & Y \arrow[rr, "\mathrm{pr}_j"] \arrow[d, "\alpha\ \in W"'] & & Y_j \arrow[d, "\beta_j\ \in \underline{W}"] \\
f, g : X \arrow[r, bend left=15] \arrow[r, bend right=15] & Y' \arrow[rr, "\varphi_j"'] & & Y'_j
\end{tikzcd}

\note{à droite de $Y'_j$, une égalité biffée, peut-être
\struck{$= \coprod_Y$} \ill{}}
On suppose que pour tout $j$, $\exists$ $\beta_j : Y_j \to Y'_j$,
$\varphi_j$, tels que $\beta_j \mathrm{pr}_j = \varphi_j \alpha$, $\beta_j
\in \underline{W}$, $\varphi_j f \underset{h}{\simeq} \varphi_j g$,
\struck{Considérons} Donc \struck{\ill{}} posant
\[
  \beta = \prod \beta_j : Y = \prod Y_j \to Y'' = \prod Y'_j
\]
\struck{\ill{}} on trouve

\page{33}

diag.\ commutatif

\begin{tikzcd}
 & Y \arrow[d, "\alpha\ \in W"'] \arrow[dr, "\beta = \prod \beta_j\ \in W"] & \\
X \arrow[r, bend left=15, "f"] \arrow[r, bend right=15, "g"'] & Y' \arrow[r, "\varphi = (\varphi_j)"'] & Y'' = \prod Y'_j
\end{tikzcd}

\note{au-dessus de la flèche $\varphi$, un mot ondulé, peut-être « comm. »}
d'où $\varphi f \simeq \varphi g$, donc $\gamma(f) = \gamma(g)$.

\emph{NB} Cela marcherait pour les produits infinis, si on savait qu'un
produit infini d'éléments de $\underline{W}_A$ est encore dans
$\underline{W}_A$.

Venons-en au lemme plus haut

(ii) $\Longleftrightarrow$ (ii bis) trivial ($\Longleftrightarrow$ si
$f \in \underline{W}_A$, alors $\forall a \in \mathrm{Ob}\,A$, $f \times
\mathrm{id}_a : X \times a \to Y \times a$ $\in \underline{W}_A$)

$\Downarrow$ (iii) on vient de le prouver
\marginal{$\Uparrow$ est évident quand on sait que $\underline{W}_A$ est
fortement saturé (conséquence de fortement saturé)}\note{en oblique dans
la marge gauche, avec une flèche $\Uparrow$ tracée en face}

\struck{Examinons (ii)} (ii bis) \struck{signifie} \uncertain{aussi} que si
$f \in \underline{W}_A$, alors $\forall Z$, $f \times \mathrm{id}_Z : X
\times Z \to Y \times Z$ est $\in \underline{W}_A$ loc.\ sur $Z$, et pour
ceci il suffit de le voir pour $Z$ dans $A$, i.e.\ $X \times a \to Y \times
a$ $\in W$.
\marginal{raisonnement déjà fait plus haut}

\begin{tikzcd}
X \times a \arrow[r] \arrow[d] & Y \times a \arrow[d] \\
X \arrow[r, "\in \underline{W}_A"'] & Y
\end{tikzcd}

Cela marche si on sait que $X \times a \to X$ et $Y \times a \to Y$
$\in \underline{W}_A$, ce qui signifie aussi que \struck{$A^{\wedge}$}
$a$ est \uncertain{tot.}\ asph.\ \struck{\ill{}}. Donc (i)
$\Longrightarrow$ (ii bis). \struck{Inversement (ii bis) implique}
\struck{En fait} Inversement, on voit de suite que si $A$ asph., alors
(ii bis) $\Longrightarrow$ (i) [prendre $f : a \to e_{A^{\wedge}}$).
donc (i) $\Longleftrightarrow$ (ii) + $A$ asphérique.

\page{34}

Reste à examiner (iv). Comme on sait déjà que (Cat) $\to
\mathrm{Hot}(\underline{W})$ commute aux produits finis, (iv) signifie
que pour $X, Y$ dans $A^{\wedge}$,
\[
  A_{/X \times Y} \to A_{/X} \times A_{/Y}
\]
\note{au-dessus, une première écriture biffée, $A_{/X} \times A_{/Y}$}
est $\in \underline{W}$. \struck{Cela implique} \struck{Dans le cas}
$X = a$, $Y = b$ dans $\mathrm{Ob}\,A$, cela signifie que $a \times b$
asphériques -- i.e.\ $A^{\wedge}$ \uncertain{tot.}\
$\underline{W}$-asph. Je dis que cela suffit, de façon précise que
$A_{/X \times Y} \to A_{/X} \times A_{/Y}$ est asph. En effet, pour un
objet $\alpha = (a_{/X}, b_{/Y})$ dans \uncertain{le second membre}, on a
\[
  (A_{/X \times Y})_{/\alpha} \simeq A_{/a \times b}
\]
\uncertain{et} on \uncertain{conclut}. En résumé

\begin{tikzcd}[column sep=small, row sep=small]
x \arrow[r] & X \times Y \\
x \arrow[r] & X \\
x \arrow[r] & Y
\end{tikzcd}

\note{petit schéma dans la marge gauche, entouré d'un trait qui porte
« si $A$ » et « $\underline{W}$-asph. », relié par une flèche à la liste
qui suit}
\begin{itemize}
  \item[(i)] $A$ \uncertain{tot.}\ $\underline{W}$-asph.
  \item[(iv)] $A^{\wedge} \to \mathrm{Hot}(\underline{W})$ commute aux
  produits binaires
  \item[(iv bis)] $A^{\wedge} \to \mathrm{Hot}(\underline{W})$ commute aux
  produits finis
\end{itemize}
$\Uparrow \qquad \Downarrow$
\begin{itemize}
  \item[(ii)] \struck{\ill{}} $\underline{W}_A$ \uncertain{stable} par
  produits finis
  \item[(ii bis)] $f \in \underline{W}_A \Rightarrow f \times \mathrm{id}_X
  \in \underline{W}_A$ ($\forall X \in \mathrm{Ob}\,\widehat{A}$)
  \item[(iii)] $A^{\wedge} \to \mathrm{Hot}_A$ commute aux produits finis
\end{itemize}
\note{les équivalences entre (i), (iv), (iv bis), puis entre (ii), (ii
bis), (iii), sont marquées de flèches doubles $\Updownarrow$ dans
l'accolade ; la flèche $\Uparrow$ de (iii) vers (ii bis) porte un mot
oblique, peut-être « \uncertain{si $W$ saturé} »}

\page{35}

\struck{Mais il n'est pas clair pour moi si}

\struck{\ill{}} \emph{NB} La condition (iv) est visiblement stable par
sommes (i.e.\ si \struck{des} $A_\alpha$ y satisfont, $\coprod_\alpha
A_\alpha$ aussi), mais la condition (i) non, donc
\struck{(i) aussi}\note{« aussi » est lu sous la rature}
l'implication (iv) $\Rightarrow$ (i) est stricte. Mais une $A$ qui
satisfait (ii), et qui est connexe, \struck{\ill{}} est-elle
$\underline{W}$-asphérique (donc \uncertain{tot.}\
$\underline{W}$-asphérique) ?

\emph{Non}\note{doublement souligné} Prenons p.ex.\ \struck{\ill{}} $A$ un
groupoïde, dans lequel $A^{\wedge}$ $\underline{W}_A$ est formé des seuls
iso (car ici $\underline{W}$ est l'équiv.\ faible \uncertain{habituelle}),
donc stable par produits finis …

\struck{L'}argument de commutation aux produits finis\note{« aux produits »
est ajouté au-dessus de la ligne} marcherait directement dans (Cat)
également, \struck{puisqu'on sait} \struck{\ill{}} qu'un produit
\struck{\ill{}}fini de foncteurs $\in \underline{W}$ est $\in
\underline{W}$. Pour le cas des produits infinis, il faut voir qu'un
produit d'éléments de $\underline{W}$ est dans $\underline{W}$, \struck{et}
et/ou itou dans un $\widehat{A}$ (le deuxième impliquera le premier).
Déjà pour l'équiv.\ faible habituelle $\underline{W}_0$, ce n'est pas
évident. Il y a des

\page{36}

chances qu'on arrive à le voir en \uncertain{prouvant} qu'un produit
$\prod_{j \in J} X_j$ d'\uncertain{objets} $\underline{W}$-asphériques est
$\underline{W}$-asphérique ? Ça me semble moins important de toutes façons
que la question des sommes infinies.

\page{38}

\ill{} \ill{} \ill{}\note{ligne d'en-tête, de la même encre, non lue ; le
dernier mot commence par « Com »}

\emph{Thm} Soient $I$ un ordonné \struck{filtrant}, $X, Y$ dans (Cat),
$(X_i)_{i \in I}$ et $Y = (Y_i)_{i \in I}$ deux filtrations par des
sous-cat.\ fermées (resp.\ ouvertes), \struck{$f$} $X = \bigcup X_i$,
$Y = \bigcup Y_i$,
\[
  f : X \to Y
\]
tel que $f(X_i) \subset Y_i$, soit $f_i : X_i \to Y_i$. Supposons que les
$f_i$ soient des équiv.\ faibles, alors $f$ est une équiv.\ faible.
\marginal{Il suffit que dans $I$ les Inf existent (\uncertain{condition}
\ill{}) et \uncertain{vérifient} que si $i = \mathrm{Inf}\, j_\alpha$, on
ait $X_i = \bigcap_\alpha X_{j_\alpha}$}\note{en oblique dans la marge
gauche, en face des six premières lignes ; « ouvertes » dans l'énoncé est
lu sans certitude}

On est ramené par dualité au cas des sous-catégories ouvertes. Dans ce
\struck{\ill{}} cas, interpréter la filtration de $X$ comme la donnée de
\[
  X \xrightarrow{\;\varphi\;} I
\]
défini par
\[
  \varphi(x) = \mathop{\mathrm{Inf}}_{i \in I(x)} i ,
  \qquad I(x) = \{\, i \in I \mid x \in X_i \,\}
\]
\note{sous « Inf », un mot biffé devant $I(x)$, et « pour » biffé dans
l'accolade}
\marginal{\emph{NB} $\varphi(x)$ est \uncertain{aussi} le plus petit
élément de $I(x)$ et $I(x) = \{\, j \in I \mid j \geq \varphi(x) \,\}$
(\uncertain{dans} \uncertain{filtration} ouvertes)}\note{à droite de la
formule ; la dernière parenthèse est reliée par un trait à « les Inf »
de la ligne suivante}
-- on suppose donc que dans $I$ les Inf existent. On a\note{« a » récrit
sur « posons » biffé}, pour
\[
  \alpha : x \to y ,
\]
$\varphi(x) \leq \varphi(y)$ (car $I(y) \subset I(x)$, donc
$\mathrm{Inf}\, I(y) \geq \mathrm{Inf}\, I(x)$), d'où la valeur de
$\varphi(\alpha)$. On \struck{aura bien}\note{« a donc » et « signifie »
sont récrits au-dessus des mots biffés, lecture incertaine}
\[
  X_i = \varphi^{-1}(I_i) \qquad \text{où } I_i = \{\, j \in I \mid j
  \leq i \,\}
\]
\struck{$\varphi$} i.e.\ $x \in X_i \Longleftrightarrow \varphi(x) \leq i$
(i.e.\ $i \in I(x)$).

\page{39}
\note{feuillet double, écrit en deux colonnes ; la colonne de gauche
achève l'argument de la p.~38, celle de droite ouvre un nouvel énoncé}

\struck{Bien sûr, par définition}

Il \uncertain{faudra} \emph{expliciter} que pour une catégorie $I$
d'indices telle que les Inf de parties non vides existent, les foncteurs
$X \to I$ sont en corr.\ 1-1 avec les filtrations, \struck{\ill{}}
\uncertain{explicites} croissantes
\[
  I \to \mathrm{Ouv}(X), \qquad i \mapsto X_i
\]
\note{devant $\mathrm{Ouv}(X)$, un premier mot biffé, peut-être
« Sscat »}
qui commutent aux dits Inf, et telles que l'on ait $X = \bigcup X_i$.
\marginal{cette dernière condition équivaut à : $\varinjlim X_i \simeq
X$}

Ceci posé, on a donc
\[
  X_{/i} = X_i
\]
\struck{d'où les restrictions} \struck{\ill{}} \struck{de} $f$
\[
  f_i : X_i \to Y_i
\]
s'identifient :
\[
  f_{/i} : X_{/i} \to Y_{/i}
\]
donc on \ill{} par L 3.

\emph{Corollaire !} Idem dans un $A^{\wedge}$, en utilisant le fait que
$i_A$ commute aux $\varinjlim$ quelc.\ et aussi aux produits fibrés.
\note{« fibrés » est écrit une seconde fois sous la ligne, à gauche ; un
trait vertical marque le corollaire en marge}

\emph{Cor ?} Si les $X_i$ sont $\underline{W}$-asphériques et $I$
connexe, alors $X$ est $\underline{W}$-asphérique (On prend $Y_i =
e_{(\mathrm{Cat})}$).

\emph{Thm} L'ens.\ $\underline{W}$ est saturé.

\emph{Cor} Pour tt $A$, $\underline{W}_A$ est saturé.

\note{la démonstration qui suit est barrée de deux longs traits obliques}
\struck{\emph{Dém}} Soit $f : X \to Y$ tel que $\gamma(f)$ iso,
\uncertain{i.e.}\ $\exists$

\begin{tikzcd}
X \arrow[r, "f"] \arrow[d, "\alpha \in \underline{W}"'] & Y \arrow[dl, "g"] \arrow[d, "\beta \in \underline{W}"] \\
X' \arrow[r, "f'"'] & Y'
\end{tikzcd}

\note{sur la page, deux croquis séparés : $X \xrightarrow{f} Y
\xrightarrow{g} X'$ avec $\alpha : X \to X'$, et $X' \xrightarrow{f'} Y'$
avec $\beta : Y \to Y'$ ; les deux sont réunis ici en un seul carré}
\struck{avec $\gamma(gf) = \gamma(\alpha)$, et}
avec
\[
  \gamma(gf) = \gamma(\alpha), \quad \gamma(f'g) = \gamma(\beta), \quad
  \beta f = f' \alpha .
\]
Quitte à changer $\alpha$ en un composé avec $X' \to X'_1$, et de même
$\beta$ en $Y'$ avec $Y' \to Y'_1$ (donc $gf \in \underline{W}$, $f'g \in
\underline{W}$), OPS $gf \sim \alpha$, $f'g \sim \beta$,
\struck{et \ill{} (quitte à changer)} \struck{\ill{} $\alpha$ et $\beta$
par $gf$ et $f'g$) \ill{}} \struck{\ill{}}

\begin{tikzcd}
X \arrow[r, "f"] \arrow[d, "\alpha"'] & Y \arrow[dl, "g"'] \arrow[d, "\beta"] \\
X' \arrow[r, "f'"'] & Y'
\end{tikzcd}

\[
  \alpha, \beta \in \underline{W}, \qquad gf = \alpha, \quad f'g = \beta
  \quad \text{(quitte à factoriser)}
\]
\marginal{$X \xrightarrow{f} Y \xrightarrow{g} Y' \xrightarrow{f'} X'$,
$gf \in \underline{W}$, $f'g \in \underline{W}$}\note{à gauche du carré ;
l'ordre $Y'$, $X'$ est celui de la page}
On pourrait d'ailleurs supposer que $f$ soit immersion ouverte, puis
(quitte à factoriser \struck{$g$, puis $f'$} \ill{}) que $f, g, f'$ sont
des immersions ouvertes.

\begin{tikzcd}[column sep=small]
X \arrow[rr, hook] & & Y \arrow[dl, "g_1"'] \\
 & X'_1 \arrow[dl] \arrow[dr] & \\
X' \arrow[rr] & & Y'
\end{tikzcd}

\note{croquis en bas de la colonne, lu sans certitude : de $X'_1$ partent
une flèche vers $X'$ et une flèche (repassée) vers $Y'$, et $X'$ va à
$Y'$}

\page{40}

Soit $f = f_0 \in \mathrm{Fl}(\mathrm{Cat})$, $f_0 : X_0 \to X_1$\note{« $:
X_0 \to X_1$ » est ajouté au-dessus de la ligne, sur un « $X =$ »
biffé}, telle que $\gamma(f)$ iso : prouver $f \in \underline{W}$. OPS $f =
f_0$ immersion ouverte. Par calcul des fractions, $\exists$ $g = g_1 :
X_1 \to X_2$, telle que $g_1 f_0 \in \underline{W}$ ; donc $\gamma(g_1
f_0) = \gamma(g_1)\gamma(f_0)$ iso, donc $\gamma(g_1)$ iso. OPS (quitte à
factoriser) que $g$ immersion ouverte. Donc en procédant de proche en
proche
\marginal{$X_0 = X \xrightarrow{f = f_0} Y = X_1 \xrightarrow{g = f_1}
X_2$}\note{petit schéma dans la marge gauche, les égalités écrites
verticalement}

\struck{$f_0, g_0, f_1, g_1, \ldots, f_n, g_n$}\note{au-dessus, une
seconde liste biffée, $f_0, f_1, f_2$}
\[
  X_0 \xrightarrow{f_0 = f} X_1 \xrightarrow{f_1 = g} X_2
  \xrightarrow{f_2} \cdots
\]
\struck{donc} système inductif d'imm.\ ouvertes, telles que les produits
de deux consécutifs soient $\in \underline{W}$. Soit $X_\infty =
\varinjlim_n X_n$, on a

\begin{tikzcd}
X_0 \arrow[rr, "f = f_0"] \arrow[dr, "i_0"'] & & X_1 \arrow[dl, "i_1"] \\
 & X_\infty &
\end{tikzcd}

commutatif ; il suffit de prouver que $i_0$ et $i_1$ sont $\in
\underline{W}$. [Il suffit] bien sûr \ill{} \struck{\ill{}}

\struck{On prouvera} \emph{On prouvera} pour $i_0$ [idem $i_1$, en
remplaçant $f_0$ par $f_1$]. On a
\[
  X_0 = \varinjlim (X_0 \to X_0 \to X_0 \to \cdots X_0 \to \cdots)
\]
et aussi
\[
  X_\infty = \varinjlim_n (X_0 \to X_2 \to X_4 \to \cdots X_{2n} \to
  \cdots)
\]
et $i_0$ est $\varinjlim$ des $i_0^{2n} : X_0 \to X_{2n}$, qui sont
$\in \underline{W}$, donc on gagne !!
\note{devant « et $i_0$ », un mot d'une autre encre, biffé et illisible}

\end{document}
